


\paragraph{Event-level precision, recall, and F1.}
A detected interval is counted as a true positive if it overlaps at least one
unmatched ground-truth blink region; unmatched detections are false positives,
and unmatched ground-truth regions are false negatives.
Per-session precision, recall, and F1 are computed from these event-level counts.


We adopt \textbf{macro-averaged F1} as the primary reporting metric.
Macro F1 computes precision, recall, and F1 independently for each session and
then reports the unweighted mean of the session-level F1 scores.
This gives every session equal weight regardless of how many blink events it
contains. The goal of epoch-based blink detection in driving studies is to produce
a detector that works reliably session by session, not one that merely accumulates
favourable aggregate counts. Micro F1, which pools all true positives, false
positives, and false negatives across sessions before computing a single ratio,
is dominated by sessions with large blink counts and can mask consistent failure
modes on sessions with atypically few blinks.


Detection was scored at the event level: a detected interval that overlapped an unmatched ground-truth interval was counted as a true positive, unmatched detections were counted as false positives, and unmatched ground-truth intervals were counted as false negatives. Event-level precision, recall, and $F_1$ were macro-averaged across sessions, and event-level $F_1$ was preferred over accuracy because blink events were sparse relative to non-blink samples. A detection was matched to a ground-truth interval at an intersection-over-union tolerance of 0.1. Methods were compared using paired two-tailed Wilcoxon signed-rank tests on session-level $F_1$, with Bonferroni correction applied to the six comparisons among the four methods and to the six non-reference durations of the epoch sweep against the 30 s reference; rank-biserial effect sizes were also reported.


Every condition is summarised at its \emph{best-channel-per-session} operating point: for
each session the channel or channel group with the highest event-level $F_1$ is selected,
and those per-session values are then averaged. The same rule is applied to all four
conditions, so no condition is given a selection advantage the others do not receive.